\documentclass[11pt,a4paper]{report}
\usepackage[T1]{fontenc}
\usepackage[utf8]{inputenc}
\usepackage[english]{babel}
\usepackage{lmodern,microtype}
\usepackage[a4paper,margin=2.25cm]{geometry}
\usepackage{amsmath,amssymb,booktabs,tabularx,longtable,array}
\usepackage{xcolor,hyperref,enumitem,fancyhdr,titlesec,tcolorbox}
\hypersetup{colorlinks=true,linkcolor=blue!40!black,urlcolor=blue!40!black}
\definecolor{deepblue}{HTML}{183B56}
\definecolor{softblue}{HTML}{EAF2F8}
\pagestyle{fancy}\fancyhf{}\lhead{ED-POMDP Research Edition v1.4}\rhead{Research Programme 2026--2030}\cfoot{\thepage}
\setlength{\parindent}{0pt}\setlength{\parskip}{0.55em}\setlength{\emergencystretch}{2em}
\setlist{leftmargin=1.5em,itemsep=0.25em,topsep=0.25em}
\titleformat{\chapter}{\Huge\bfseries\color{deepblue}}{\thechapter}{0.8em}{}
\titleformat{\section}{\large\bfseries\color{deepblue}}{\thesection}{0.65em}{}

\title{\textbf{ED-POMDP Research Programme 2026--2030}\\
\large Revised gated roadmap after the Step 2 benchmark}
\author{Guillaume Harbonnier -- STRAT-Q Research Programme}
\date{2 August 2026}

\begin{document}
\maketitle
\tableofcontents

\chapter{Programme status}

Steps 1 and 1.1 are shipped. Step 2 is closed after frozen execution, independent review, mechanism diagnosis, and claim adjudication. Step 2 did not validate broad ED-POMDP superiority.

\begin{tcolorbox}[colback=softblue,colframe=deepblue,boxrule=0.8pt,arc=2mm]
\textbf{Programme gate.} The former direct path from synthetic benchmark to industrial calibration is suspended. Milestone 2R---Theory and Claim Reset---is mandatory before any new confirmatory benchmark or industrial milestone.
\end{tcolorbox}

The benchmark contained 3,360 raw policy episodes:
\[
4\ \mathrm{regimes}\times 4\ \mathrm{budgets}\times 30\ \mathrm{seeds}\times 7\ \mathrm{policies}=3{,}360.
\]
The mechanism analysis contained 2,400 ED-POMDP-versus-baseline pairs:
\[
4\times4\times30\times5=2{,}400.
\]
Of 1,119 Brier-improving pairs, 1,054 retained the same terminal action.

\chapter{Scientific posture}

The programme preserves the distinction between latent system risk and latent evidence-production quality. It withdraws the assumption that improved evidence selection or calibration naturally creates better decisions under an unchanged terminal architecture.

The revised central question is:
\begin{quote}
Under which conditions does additional or better-qualified evidence change a governed engineering decision, and does that changed decision reduce materially important loss at comparable total cost?
\end{quote}

\chapter{Milestone 2R -- Theory and Claim Reset}

\section{Objectives}

\begin{enumerate}
  \item Reconstruct the complete causal chain from evidence acquisition to realised loss.
  \item Audit thresholds, loss ratios, stopping, evidence costs, dependence, redundancy, and decision admissibility.
  \item Replace broad superiority claims with conditional hypotheses H3-A to H3-E.
  \item Define new progression gates and seed-separation rules.
\end{enumerate}

\section{Exit criteria}

Milestone 2R exits only when:
\begin{itemize}
  \item the P0--P5 bet register is approved;
  \item structural assumptions and boundary reachability are audited;
  \item candidate hypotheses have explicit estimands and refutation criteria;
  \item Step 2 headline seeds are cryptographically and procedurally excluded from tuning;
  \item a Step 3A exploratory protocol is independently reviewable.
\end{itemize}

Milestone 2R produces no new confirmatory performance claim.

\chapter{Post-reset work packages}

\section{WP3A -- Exploratory mechanism benchmark}

\textbf{Purpose:} determine when belief changes cross decision boundaries, alter stopping, trigger compensating controls, or reduce loss.

\textbf{Design:} new development seeds; adaptive stopping; heterogeneous costs; correlated and redundant evidence; joint risk/evidence-quality terminal rules; explicit CONDITIONAL GO controls.

\textbf{Constraint:} no confirmatory claim.

\section{WP3B -- New preregistered benchmark}

\textbf{Primary endpoints:}
\begin{enumerate}
  \item weighted terminal decision loss;
  \item unsafe-GO rate;
  \item total evidence cost.
\end{enumerate}

Brier score and ECE are secondary mechanism diagnostics. Confirmatory seeds remain untouched until protocol and analysis freeze.

\section{WP4 -- Realistic STRAT-Q scenarios}

The realistic benchmark adds heterogeneous evidence costs, correlated tests and OTEL signals, representative environments, governed provenance, schedule constraints, compensating controls, and project-specific loss models.

\section{WP5 -- Industrial replay and shadow evaluation}

This work requires a passed data-readiness gate and prior bounded validation. The industrial pilot evaluates an architecture that has already survived synthetic testing; it must not be used to rescue an unstable theory.

\chapter{Candidate hypotheses}

\begin{longtable}{p{0.13\textwidth}p{0.75\textwidth}}
\toprule
ID & Candidate hypothesis \\
\midrule
H3-A & A terminal rule using both system-risk and evidence-quality belief reduces weighted loss under evidence degradation relative to a risk-only rule. \\
H3-B & Adaptive stopping reduces mean evidence cost without increasing terminal loss or unsafe-GO rate. \\
H3-C & CONDITIONAL GO linked to explicit, verifiable compensating controls reduces loss relative to an unstructured intermediate decision region. \\
H3-D & VOI provides benefit only when evidence is heterogeneous, discriminating, non-redundant, and boundary-capable. \\
H3-E & Evidence-quality-aware terminal rules reduce unsafe-GO rate under controlled degradation at comparable total cost. \\
\bottomrule
\end{longtable}

\chapter{Epistemic governance}

Evidence level records source class, not support. Step 2 changes the adjudicating evidence level from \texttt{NONE} to \texttt{SYNTHETIC} because a frozen synthetic benchmark now exists. Polarity and disposition remain separate. A claim can therefore have synthetic evidence and remain unsupported.

The canonical combinations are:
\begin{itemize}
  \item \texttt{CLM-VOI-001}: synthetic evidence, adverse/mixed polarity, \texttt{NOT\_SUPPORTED\_STEP2};
  \item \texttt{CLM-EQ-001}: synthetic evidence, mixed narrow-positive polarity, \texttt{BROAD\_FORM\_NOT\_SUPPORTED\_NARROW\_ECE\_SIGNAL}.
\end{itemize}

\chapter{Anti-leakage and reproducibility}

\begin{itemize}
  \item Step 2 headline seeds remain immutable historical evidence.
  \item Step 3A uses new development seeds.
  \item Step 3B uses untouched confirmatory seeds.
  \item Protocol, code, endpoints, multiplicity, and stopping criteria freeze before confirmation.
  \item Simple governed baselines remain first-class comparators.
  \item A new negative result is a valid programme outcome.
\end{itemize}

\chapter{Milestones}

\begin{longtable}{p{0.13\textwidth}p{0.25\textwidth}p{0.49\textwidth}}
\toprule
ID & Milestone & Acceptance evidence \\
\midrule
M1 & Scientific foundations & Shipped. \\
M1.1 & Identifiability foundations & Shipped. \\
M2 & Frozen matched-budget benchmark & Closed; broad superiority claims not supported. \\
M2R & Theory and claim reset & Audited assumptions, P0--P5 register, H3-A--H3-E candidates, progression gates. \\
M3A & Exploratory mechanism benchmark & New development seeds; no confirmatory claim. \\
M3B & Preregistered confirmatory benchmark & Untouched seeds; primary loss, safety, and cost family. \\
M4 & Realistic STRAT-Q scenario & Governed heterogeneous evidence and project-specific loss. \\
M5 & Industrial replay and shadow mode & Passed data gate and prior bounded validation. \\
\bottomrule
\end{longtable}

\chapter{Conclusion}

Step 2 performed its scientific function by exposing the proposed advantages to a frozen attempt at falsification. The programme continues, but not on the former assumptions. Milestone 2R makes the negative result operationally useful by redesigning the complete belief-to-action architecture before any new claim or deployment path.

\end{document}
